\section{Realization in a polar product}
\label{sec:realization}

We now combine the ball construction with the uniform slice estimate.
Writing \(z=u+ix\), the imaginary coordinates \(x\) will give the position
in a convex body.  The signed slice integrals will correct the standard
momentum \(-u\) to account for the holomorphic part of the symplectic form.

\begin{proposition}\label{prop:finite-strips}
Let $n\ge2$ be an integer, let $b_1,\ldots,b_m\in\R^n$ span $\R^n$,
and set
\[
 K=\{x\in\R^n: |b_j\cdot x|\le 1,
                         \quad 1\le j\le m\}.
\]
For every $0<c<4$ there is a smooth symplectic embedding
\[
 B^{2n}(c)\longrightarrow
 \interior(K)\times\interior(K^\circ).
\]
\end{proposition}

\begin{proof}
\noindent\textit{Step 1: encode the strips by holomorphic functions.}
The body $K$ and its finite list of defining vectors are fixed throughout
the proof. Extend each real functional $b_j\cdot x$ complex-linearly to
$\ell_j(z)=b_j\cdot z$. Let $F:\D\to D$ and $g=F^{-1}$ be the maps of
Lemma~\ref{lem:planar-map}, and define
\begin{equation}\label{eq:strip-domain}
 \Omega=\{z\in\C^n:\ell_j(z)\in D\text{ for every }j\}.
\end{equation}
This is an open neighborhood of the origin. The spanning condition makes
the linear map $z\mapsto(\ell_1(z),\ldots,\ell_m(z))$ injective, so there
is a constant $C_b$ such that
\[
 |z|\le C_b\max_j|\ell_j(z)|.
\]
Since $D$ is bounded, $\Omega$ is bounded as well; neither set depends on
the integer $k$ chosen below. The same linear bound makes $K$ compact,
and a sufficiently small Euclidean ball lies in all its defining strips.
In particular, $0\in\interior(K)$.

For each integer $k\ge2$, put
\[
 f_j(z)=g(\ell_j(z))^k,\qquad
 \tau(z)=\sum_{j=1}^m|f_j(z)|^2,\qquad
 V=\{z\in\Omega:\tau(z)<1\},\qquad
 \omega=\ddc\bigl(|z|^2+\tau\bigr).
\]
The tuple $f=(f_1,\ldots,f_m)$ is holomorphic on $\Omega$. Since $g$ has
its only zero at the origin, the common zero set of $f$ is
\[
 \{z\in\Omega:\ell_j(z)=0\text{ for every }j\}=\{0\}.
\]
Holomorphicity at the origin gives $|f(z)|=O(|z|^k)$. To verify the
compactness hypothesis of Theorem~\ref{thm:holomorphic-ball}, fix
$0<s<1$. On $\{\tau\le s\}$, for every $j$ we have
\[
 \ell_j(z)\in
 F\bigl(\{w\in\C:|w|\le s^{1/(2k)}\}\bigr),
\]
a compact subset of $D$. The linear bound above makes this sublevel
bounded. Moreover, the limit of any convergent sequence in it still has
all its images under the functionals $\ell_j$ in that compact subset
of $D$, hence belongs to
$\Omega$ and satisfies $\tau\le s$ by continuity. Thus the sublevel is
compact in $\Omega$. Theorem~\ref{thm:holomorphic-ball} now supplies
symplectic embeddings into $(V,\omega)$ of every standard ball of
capacity less than $\pi k$.

\medskip
\noindent\textit{Step 2: realize the form by a globally injective map.}
Write $z=u+ix$ and use the signed integrals $J_k$ of
Proposition~\ref{prop:uniform-slice} to define
\begin{equation}\label{eq:strip-phase-map}
 \Phi_k:V\longrightarrow\R^n_q\times\R^n_p,
 \qquad
 q=x,\qquad
 p=-u-\sum_{j=1}^m
       J_k(b_j\cdot u,b_j\cdot x)b_j.
\end{equation}
The signed integrals $J_k$ are smooth on $D$ by
Proposition~\ref{prop:uniform-slice}, so $\Phi_k$ is smooth.

Set $v_j=b_j\cdot u$ and $t_j=b_j\cdot x$. Pulling back the target
form gives
\begin{align}
 \Phi_k^*\omega_0
 &=\sum_{\ell=1}^n dx_\ell\wedge(-du_\ell)
       -\sum_{j=1}^m dt_j\wedge dJ_k(v_j,t_j)\notag\\
 &=\sum_{\ell=1}^n du_\ell\wedge dx_\ell
       +\sum_{j=1}^m
          \partial_vJ_k(v_j,t_j)\,dv_j\wedge dt_j\notag\\
 &=\ddc\bigl(|z|^2+\tau\bigr)=\omega.
 \label{eq:strip-pullback}
\end{align}
Indeed, the terms containing $\partial_tJ_k$ vanish because
$dt_j\wedge dt_j=0$, while
\[
 \partial_vJ_k(v,t)=k^2|g(v+it)|^{2k-2}|g'(v+it)|^2.
\]
This is precisely the coefficient of $\ddc|g(v+it)^k|^2$ in front of
$dv\wedge dt$, with the convention for $\dc$ fixed above.

The map $\Phi_k$ is globally injective. To see this, equality of its
$q$ coordinates first fixes $x$. For two possible real coordinates
$u_1,u_2$ at that $x$, write
\[
 \Delta u=u_2-u_1,\qquad
 v_{ji}=b_j\cdot u_i,\qquad t_j=b_j\cdot x,
\]
and denote the corresponding momenta by $p_1,p_2$. Both $v_{j1}$ and
$v_{j2}$ belong to the same horizontal interval of $D$ at height $t_j$.
Since $\partial_vJ_k\ge0$ on this entire interval,
\begin{align*}
 (p_2-p_1)\cdot\Delta u
 &=-|\Delta u|^2
   -\sum_{j=1}^m
     \bigl(J_k(v_{j2},t_j)-J_k(v_{j1},t_j)\bigr)
     (v_{j2}-v_{j1})\\
 &\le -|\Delta u|^2.
\end{align*}
Thus equal momenta force $u_1=u_2$. This comparison applies to any two
points of the fiber, even if that fiber of $V$ is disconnected.
Equation~\eqref{eq:strip-pullback} and nondegeneracy of $\omega$ show that
$\Phi_k$ is a local diffeomorphism. It is therefore open, and its
injectivity gives a smooth inverse onto its image. Consequently
$\Phi_k$ is a symplectic embedding.

\medskip
\noindent\textit{Step 3: control the momentum by the polar body.}
We next bound its image. All imaginary parts of points of $D$ lie in
$(-1,1)$, so $|b_j\cdot x|<1$ for every $j$. The finite system of
strict inequalities implies $q=x\in\interior(K)$. Introduce the support
function
\[
 h_K(p)=\max_{y\in K}\ip{p}{y}.
\]
It satisfies $h_K(\pm b_j)\le1$, and
$\interior(K^\circ)=\{p:h_K(p)<1\}$. Because $\Omega$ is bounded,
there is a finite constant $C_K$, independent of $k$, such that
$h_K(-u)\le C_K$ on $\Omega$. Subadditivity and positive homogeneity of
$h_K$, followed by Proposition~\ref{prop:uniform-slice}, give
\begin{align}
 \frac{h_K(p)}{k}
 &\le \frac{C_K}{k}
       +\sum_{j=1}^m\frac{|J_k(v_j,t_j)|}{k}\notag\\
 &\le \frac{C_K}{k}
       +\frac{\pi}{4}\sum_{j=1}^m|g(\ell_j(z))|^{2k}
       +m\epsilon_k
  =\frac{C_K}{k}+\frac{\pi}{4}\tau(z)+m\epsilon_k.
 \label{eq:strip-support}
\end{align}
Here $\epsilon_k\ge0$ tends to zero, uniformly in the planar argument.
The number $m$ and the constant $C_K$ are fixed before $k$ varies.
It follows that, for any $\eta>0$ and all sufficiently large $k$,
\[
 \frac{C_K}{k}+m\epsilon_k<\frac{\eta\pi}{4}.
\]
Since $\tau<1$ on $V$, Equation~\eqref{eq:strip-support} then gives the
strict inclusion
\begin{equation}\label{eq:strip-scaled-image}
 \Phi_k(V)\subset
 \interior(K)\times S\interior(K^\circ),
 \qquad S=(1+\eta)\frac{\pi k}{4}.
\end{equation}

\medskip
\noindent\textit{Step 4: rescale the ball and the momentum together.}
Finally, fix $0<c<4$. Choose $\eta>0$ with $(1+\eta)c<4$, and then fix
an integer $k$ large enough for Equation~\eqref{eq:strip-scaled-image}.
The resulting constant $S$ satisfies $Sc<\pi k$, so there is a
symplectic embedding
\[
 E:B^{2n}(Sc)\longrightarrow(V,\omega).
\]
Let $D_{\sqrt S}(z)=\sqrt S\,z$ and
$L_S(q,p)=(q,p/S)$. The composition
\begin{equation}\label{eq:strip-rescaling}
 L_S\circ\Phi_k\circ E\circ D_{\sqrt S}:
 B^{2n}(c)\longrightarrow
 \interior(K)\times\interior(K^\circ)
\end{equation}
is a smooth embedding. Since $D_{\sqrt S}$ maps $B^{2n}(c)$ onto
$B^{2n}(Sc)$ and
\[
 D_{\sqrt S}^*\omega_0=S\omega_0,
 \qquad L_S^*\omega_0=S^{-1}\omega_0,
\]
its pullback of $\omega_0$ is exactly $\omega_0$. This proves the
proposition.
\end{proof}
